\chapter{动力学内核：统一状态算子与多尺度闭环}
\label{chp:dynamics_kernel_ch}

前两章分别给出了状态演化与状态—结构耦合，本章进一步回答一个更严格的问题：怎样把活动状态、结构参数、规范约束、外部输入和控制决策放进同一套可计算内核，同时不把智能系统误写成某种物理场的同构物。我们保留原章关于即时推理、长期学习、价值更新、主动控制以及环境闭环的全部问题，也保留其快慢尺度和双向反馈案例；但原先借用狄拉克、爱因斯坦、杨—米尔斯方程所得的四个名称，不再承担证明责任。新的内核以带类型的混合状态、任务目标、约束集合和提交协议为对象：连续方程描述固定结构区间内的近似变化，离散事件描述结构编辑、权限变化和外部行动，经验主张必须由辨识与干预验证。HSF-HD在此提供智能的一般状态动力学，AgentOS只是用特定组件实现它的一种工程路径。

\section{带类型状态空间、接口与可辨识对象}
\label{sec:section_9973}

我们先定义模型对象，而不是从物理名词反推认知含义。令运行时状态为
\begin{equation}
X(t)=\bigl(h(t),S(t),Z(t),B(t),g(t),r(t),M_e(t)\bigr)\in\mathcal X ,
\end{equation}
其中 $h\in\mathbb R^{d_h}$ 是有限工作窗口内的活动状态，$S=(V,E,\rho,\nu)$ 是带节点身份、关系类型、角色和版本的局部结构网络，$Z\in\mathbb R^{n\times d_z}$ 是全局分布表示，$B\in[0,1]^{|V|\times n}$ 是结构身份与分布候选之间的绑定，$g\in\mathcal G$ 是目标与规范状态，$r\in\mathbb R_+^q$ 是延迟、内存、调用额度等资源记录，$M_e$ 是外部记忆。各分量的坐标可以不同，但每个读出都必须声明输入版本、单位和有效域。状态范数只是选定坐标下的活动尺度，不等于概率、意义、价值或焦耳能耗。

观测写为 $o_k=(y_k,\tau_k,s_k,q_k)$，分别表示测量值、时间戳、来源和质量标记；行动写为 $a_k=(u_k,\alpha_k,\pi_k)$，分别表示执行参数、权限凭据和回滚策略。编码器 $E_o$ 把原始观测变成候选证据，读出器 $R_Q$ 针对任务查询集合 $Q$ 产生答案或行动资格：
\begin{equation}
e_k=E_o(o_k;\theta_E),\qquad \hat y_{Q,k}=R_Q(X_k,e_k).
\end{equation}
这是局部结构网络与全局分布表示的接口，而不是二者相同的声明。$S$负责身份、角色、组合和来源依赖，$Z$负责相似性、候选和不确定性，$B$使任务相关身份能够被联合读出。仅有向量相似不能证明对象同一，仅有图节点也不能提供未见样本的连续泛化。

模型中至少有四类不同时间。外部时钟 $t$ 以秒或事务序号计量；快活动步长为 $\Delta t_f$；结构更新周期为 $\Delta t_s=m\Delta t_f$；规范与策略评审按事件触发而非固定微分。若传感器、模型服务和外部记忆异步，状态还要携带各自水位与最大延迟。把它们写成一个连续“认知时间”会隐藏乱序、重试和过期证据，因而不能支持工程验证。

可辨识性先于漂亮方程。若两个状态 $X$ 与 $X'$ 对所有允许查询、干预和未来观测都给出同一分布，它们在当前任务上不可区分；此时强行解释每个坐标没有经验依据。相反，当结构编辑改变了身份或可达关系，即使短期输出相同，也应保留两个版本，因为后续干预可能揭示差异。我们用任务等价关系
\begin{equation}
X\sim_Q X'\iff \sup_{q\in Q,\,\iota\in\mathcal I}d_q\!\left(P_q(\cdot\mid X,\iota),P_q(\cdot\mid X',\iota)\right)\le\varepsilon_q
\end{equation}
声明压缩条件。这里 $\mathcal I$ 是允许的干预集，$d_q$ 是有单位或经归一化的查询距离。该定义说明所谓“几何”是任务相关区分结构；它可以由图距离、统计度量或控制可达性实现，不必假定伪黎曼流形、旋量或非阿贝尔规范群。

外部输入和宏观控制也必须分型。输入 $e_k$ 是证据，不因强度大就自动真实；控制 $c_k$ 是系统选择的干预，不因目标明确就自动有权限。在AgentOS实例中，$h(t)$、$E$、$\Theta$、$\Psi$、SPC、TCE、CCM、Boundary、Offloading与$M_e$分别承担工作状态、事件证据、稳定结构、候选生成、压缩读出、控制、复杂度限制、边界审核、卸载和外部记忆等角色。其他智能系统可以用不同模块，只要能给出等价的状态、接口、验证和失败恢复语义。

这些对象还需要量纲检查。概率和置信度是无量纲数，延迟以时间计，容量以字节或记号数计，图距离依赖边权定义，设备能耗以焦耳计。若把它们统一标准化用于优化，日志仍须保存原始量与换算函数，否则一次阈值越界无法追溯。状态空间的维度也不是能力等级：维度增加可能只代表冗余编码，节点增加可能只代表重复实体。有效容量应通过可区分查询、干预响应和泛化误差评估。

以“桌上红杯”为例，视觉编码可产生多个红色区域向量，结构网络保存杯、桌面与空间关系，绑定矩阵依据时间连续、遮挡恢复和指称历史选择对象。若杯被移动，观测坐标改变但身份可保持；若出现第二只相同杯，仅凭向量近邻就会混淆。这个案例说明结构、分布和绑定各自不可省略，也给出了通过交换位置、遮挡和改名干预来检验身份的办法。

\section{从“总作用量”到有限时域目标与约束账本}
\label{sec:section_1047}

原章试图用一个总拉格朗日量同时描述思维、结构、价值和意志。我们保留“联合权衡”的动机，但把它改写为有限时域决策问题。给定当前版本 $v_k=(S_k,g_k,\theta_k)$、预测时域 $H$ 和候选控制序列 $c_{k:k+H-1}$，定义
\begin{align}
J_k={}&\mathbb E\!\left[\sum_{j=0}^{H-1}\ell_{task}(X_{k+j},c_{k+j};v_k)+\ell_{risk}(X_{k+j},c_{k+j})\right] \notag\\
&+\ell_T(X_{k+H};v_k)+\lambda_c C(c_{k:k+H-1})+\lambda_\Delta D(v_{k+H},v_k).
\end{align}
$\ell_{task}$衡量明确任务误差，$\ell_{risk}$记录安全与外部性，$C$记录计算、调用、延迟或实测能耗，$D$约束结构和规范漂移。不同量只有在说明归一化方法后才能相加；如果某项必须严格满足，就把它写入可行集而不是用有限权重“折价”。

可行控制集合写为
\begin{equation}
\mathcal C_k=\{c:\;A_{hard}(X,c,g_k)=1,\;r(c)\le r_k,\;\alpha(c)\in\mathrm{Permit}_k\}.
\end{equation}
其中 $A_{hard}$ 表示不可违反的安全、法律或协议约束，$r(c)$ 是带单位的资源需求，权限凭据 $\alpha$ 必须通过 Boundary 验证。若 $\mathcal C_k$ 为空，正确结果是拒绝、请求补充信息或升级处理，而不是继续最小化一个把禁令变成软惩罚的标量目标。多主体情形还要维护每一项承诺的主体、对象、期限和撤销条件，不能用单一“价值场强”掩盖不可通约的义务。

目标不是静态宇宙定律。若任务版本、结构、输入分布或评价度量变化，完整导数包含
\begin{equation}
\frac{dJ}{dt}=\langle\nabla_XJ,\dot X\rangle+\partial_tJ+\langle D_SJ,\dot S\rangle+\langle D_gJ,\dot g\rangle+\langle D_oJ,\dot o\rangle .
\end{equation}
结构 $S$ 和规范 $g$ 多数时候以跳变更新，实际计算应在事件前后比较 $J(X^-,v^-)$ 与 $J(X^+,v^+)$，而不是把离散编辑伪装成普通导数。旧目标下降不能证明新任务成功；新指标提高也可能来自评价口径改变。每次优化因此保存目标版本、数据窗口、约束集合和度量定义。

联合目标还需要防止“代理指标劫持”。例如客服系统降低平均响应时间，可能通过过早结束复杂请求实现；学习系统降低预测误差，可能依赖泄漏标签；研究Agent增加引用数，可能重复引用无关来源。我们要求任务层至少包含结果准确性、过程资格和证据完整性三个读出，并用独立验证器检查。验证器不共享同一条未经确认的证据链，否则两个低误差只是共同偏差。

所谓最小作用在这里仅保留为一种规划类比：在给定时域、模型和约束下选择代价较低的轨迹。它不推出真实智能必然优化某个全局标量，也不推出目标函数具有本体地位。经验上是否近似优化，要比较候选模型对干预数据的预测；工程上采用何种目标，则是设计选择。数学结果、经验主张和设计决策由此分开。

为了避免每步重求全局计划，系统采用滚动时域：求得 $c^*_{k:k+H-1}$ 后只执行第一项，读取新观测再规划。若外部行动不可逆，执行前还要做因果影响和权限审查；内部回滚不能恢复已经发送的邮件、转账或物理动作。这个边界把“意志做功”还原为有预算、有资格且可验证的控制，而不是抽象负熵。

有限时域还暴露了终端代价的假设。若 $H$ 太短，控制器可能为了眼前收益推迟风险；若终端模型过强，又可能把未经验证的远期预测当事实。工程上应报告时域、折扣、终端集合和预测误差，并用不同 $H$ 做敏感性分析。目标之间若存在字典序，例如先保证人身安全再优化效率，求解器必须保持该顺序，而不是由权重偶然近似。

代价账本与会计账本一样需要守恒范围。一次卸载可能降低本地延迟，却增加网络费用、隐私暴露和远端能耗；一次压缩可能减少记忆，却增加重构误差。我们把成本归属到具体主体和时间窗，避免把外部性移出模型便宣称总成本下降。只有在相同边界、相同单位和相同任务版本下，两个策略的 $J_k$ 才可直接比较。

\section{快速状态算子：推理、传播与现实校正}
\label{sec:section_7654}

固定结构与规范版本 $(S_k,g_k)$ 的短区间内，我们用一般状态算子描述即时推理。令 $x(t)=(h(t),\operatorname{vec}Z(t),\operatorname{vec}B(t))\in\mathbb R^d$，假设在区间内维度不变、输入分段连续、向量场局部Lipschitz，则
\begin{equation}
\dot x=f_{\theta}(x;S_k,g_k)+G_o(x)e(t)+G_c(x)c(t),\qquad \hat y=R_Q(x,S_k,g_k).
\end{equation}
$f_\theta$是内部传播，$G_o e$是观测校正，$G_c c$是控制作用。这个方程不要求 $x$ 是波函数；若选用复值状态、振荡网络或注意力矩阵，必须另行给出坐标、归一化和读出解释。

局部结构网络上的一个具体实现是
\begin{equation}
\dot h=-Dh+\sigma\!\left(W_Sh+W_ZR_B(Z)+Ue+Vc+b\right),
\end{equation}
其中 $D\succeq0$ 是泄漏矩阵，$W_S$ 只沿结构允许的边传播，$R_B$ 按绑定读取全局候选，$\sigma$ 是有界非线性。若 $\sigma$ 的Lipschitz常数为 $L_\sigma$，并且 $\lambda_{\min}(D)>L_\sigma\|W_S\|$，无输入差分系统局部收缩；有界输入则给出有界响应。这是给定模型假设下的数学结果，不是对所有推理过程的经验定律。

现实输入不是“打破幺正性的激波”，而是带来源和噪声模型的观测。系统先计算创新量
\begin{equation}
\epsilon_k=y_k-H_k\hat x_{k|k-1},\qquad x_{k|k}=x_{k|k-1}+K_k\epsilon_k,
\end{equation}
其中 $H_k$ 是观测模型，$K_k$ 由噪声、来源质量和当前不确定性决定。异常大的残差既可能表示世界变化，也可能表示传感故障、提示注入或时间戳错位。只有通过多源核验、故障模型和后续预测，系统才决定吸收、隔离或升级该证据。输入幅度不等于真实性。

快思考与慢思考可被重写为控制预算不同的运行区间。低风险熟悉任务中，系统允许较小搜索宽度和较低验证开销；高风险或模型冲突时，CCM提高候选数，TCE分配验证预算，必要时把问题卸载给外部工具或人。所谓顿悟不是穿过拓扑虫洞，而是在候选生成、结构重组或新证据加入后出现显著更好的可检验解释。它仍需在保留集、反例和新样本上验证。

离散实现用
\begin{equation}
x_{k+1}=x_k+\Delta t_f f_\theta(x_k;S_k,g_k)+G_{o,k}e_k+G_{c,k}c_k+\xi_k .
\end{equation}
$\Delta t_f$是明确采样步长，$\xi_k$汇总数值和模型残差。连续稳定不保证任意步长稳定；在线性模态 $\dot x=-Ax$ 中，显式Euler要求每个特征值满足 $|1-\Delta t_f\lambda_i(A)|<1$。异步服务还要把最大延迟纳入扩展状态。系统保存随机种子、模型版本、工具返回摘要和重试序列，使一次“思维流”可以复查。

该快速算子保留原章的五类影响：内部联想对应 $f_\theta$，结构约束对应 $S$ 与 $W_S$，目标偏置对应 $g$，主动注意对应 $c$，感知惊奇对应 $e$。不同之处在于这些影响都有可观测接口和失败模式，不再被称作五种物理力。推理成功最终由外部任务结果、证据一致性与权限合规共同判断，而不是由状态轨迹看起来平滑来判断。

候选传播还可能发生塌缩或爆炸。若所有查询都读出同一高频候选，数值可以稳定但语义分辨率已经丢失；若循环边反复放大同一证据，活动有界也可能造成来源重复计数。系统因此监测有效候选数、来源独立数、绑定熵和循环增益，并在异常时冻结提交、降低增益或切换到显式搜索。这里的熵是候选分布的统计量，不是热力学熵。

对逻辑证明任务，快速算子还必须携带依赖图。一个结论的高激活不能替代推导，每次传播要记录使用的前提、规则和版本；若前提撤销，相关结论进入失效队列。对开放问答，证据链可能不完备，系统应输出不确定性和待核验项。这样，快速动力学既能支持联想生成，也不会把生成流畅度当作证明有效性。

如果读出同时服务多个任务，传播权重还应随查询改变。面向事实检索的状态路径不必等同于面向创意生成的路径；共享底层表示并不意味着共享同一校准阈值。系统为每类查询保存读出头、容差和验证器，防止某一任务上的高分被错误外推为一般能力。

\section{慢速结构算子：学习、记忆与可回滚编辑}
\label{sec:section_1246}

学习改变的是后续状态转移所依赖的参数和结构。令连续参数为 $\theta_s\in\Theta$，离散结构为 $S\in\mathcal S$，经验缓冲为 $E_k$。一个双时间尺度模型是
\begin{align}
x_{k+1}&=F(x_k,e_k,c_k;\theta_{s,k},S_k),\\
\theta_{s,k+1}&=\Pi_\Theta\!\left[\theta_{s,k}-\eta_k\widehat\nabla_{\theta_s}L(E_k;\theta_{s,k},S_k)+\eta_k q_k\right],
\end{align}
其中 $\eta_k$ 是无量纲学习率，$q_k$ 是稳定性、稀疏性或保留约束产生的修正，$\Pi_\Theta$ 投影到合法参数域。只有当梯度估计偏差受控、步长条件成立且数据分布近似稳定时，才可讨论随机逼近收敛；在线环境中这些条件往往只在局部窗口成立。

短时活动、参数学习和结构编辑必须分开。短时记忆可以由 $h$ 与缓存维持，不必修改长期参数；反复练习只有在独立验证表明可迁移收益时才提交到 $\theta_s$；节点、边、类型或工具契约的增加删除属于结构提案 $\Delta S$。结构提案经过影子执行：复制当前版本，迁移身份绑定，重放关键查询，然后比较
\begin{equation}
D_{keep}=\sup_{q\in Q_{keep}}d_q\bigl(R_q(X^-),R_q(X^+)\bigr),
\end{equation}
并检查新任务收益、权限和资源。只有保留差异在容差内且收益通过独立测试才提交，否则回滚。

原章的“熟能生巧、单次深刻事件、主动学习”对应三种证据制度，而不是三种压弯时空的应力。重复样本提供低方差但可能高度相关的证据；单次高风险事件可以触发快速隔离和高优先级复核，但不应仅凭情绪或激活强度永久改写模型；主动学习由信息增益、覆盖缺口和查询成本选择新证据。三者都要记录来源独立性，避免把同一内容的多次转发误当成重复确认。

稳定性—可塑性冲突可用保留损失表达：
\begin{equation}
L_{total}=L_{new}+\lambda_{keep}L_{replay}+\lambda_{id}L_{binding}+\lambda_{simp}C(S,\theta_s).
\end{equation}
$L_{new}$评价新任务，$L_{replay}$评价历史关键样本，$L_{binding}$防止身份漂移，$C$限制复杂度。权重不是自然常数，必须通过风险和验证集选择。目标下降仍不等于学习成功：系统可能记住训练样本、破坏罕见能力或利用评价漏洞，因此还需组合泛化、反事实和长期回归测试。

结构改变时不能沿用固定维度导数。新增概念意味着状态空间和绑定接口变化，系统要生成版本映射 $M_{v\to v+1}$，说明哪些节点保持身份、哪些字段被拆分、哪些查询不再可比。所谓范式转移若没有这种映射，只是不可审计的重命名。对无法无损迁移的内容，外部记忆 $M_e$ 保存旧版本和来源，使系统能够比较新旧解释而不是覆盖历史。

在AgentOS实现中，事件证据进入 $E$，稳定结构写入 $\Theta$，外部记忆承载大体量重放集，Boundary审核跨域数据，Offloading安排离线训练或人工评审。一般理论只要求存在等价的经验、结构、提交和回滚接口。生物学习、群体制度和软件Agent可采用完全不同的载体机制，不能由“都表现为慢变量”推出机制同一。

遗忘也应分型。参数衰减、索引淘汰、来源过期、检索失败和权限撤销都可能表现为“想不起来”，但修复方式完全不同。参数衰减需要重训，索引淘汰可重建索引，来源过期需要重新取证，权限撤销则不应绕过。系统在 $E$ 和 $M_e$ 中保存失效原因，避免把可恢复的访问故障误判为知识已被抹除。

学习率与更新频率还受数据相关性影响。连续接收同一事件的复制品会低估方差，使结构过快固化；延迟标签则可能把后来的结果归因给错误行动。经验缓冲应按来源簇、时间窗和干预批次抽样，提交时报告有效样本量而非原始条数。只有当回归测试覆盖罕见条件、旧版本可恢复且新收益在独立数据上成立，慢更新才算验收通过。

\section{规范更新算子：偏好、承诺与冲突恢复}
\label{sec:section_2317}

价值与目的不能用无单位的规范势表示后直接套用场方程。我们把规范状态写为 $g=(G_h,w,P,C_m)$：$G_h$ 是硬约束集合，$w$ 是可调整偏好参数，$P$ 是权限和角色，$C_m$ 是多主体承诺账本。对候选行动 $a$，资格判定先计算 $a\in\mathcal A_{legal}(X,G_h,P,C_m)$，再对剩余候选比较向量效用
\begin{equation}
U(a\mid X,g)=\bigl(u_{task},-u_{risk},-u_{cost},u_{commit}\bigr).
\end{equation}
向量可以按公开规则排序、采用Pareto筛选或进入协商，但不能在未声明尺度时压成一个“价值能量”。

可调整偏好 $w$ 的更新来自结果证据、明确反馈和规范修订，而不是“行为做久了便自动成为价值”。令 $z_k$ 是经来源核验的反馈特征，$\hat z_k(w)$ 是当前偏好模型预测，则
\begin{equation}
w_{k+1}=\Pi_W\left[w_k-\eta_g\nabla_w\ell_g(\hat z_k(w_k),z_k)-\eta_g\lambda_g(w_k-w_{ref})\right].
\end{equation}
$w_{ref}$保存经批准的参考版本，$\Pi_W$维持范围与单调性约束。行为和偏好长期不一致可能提示模型错误、外部胁迫、角色冲突或真实偏好变化；系统必须区分这些解释，不能把认知失调当作必然驱动价值同化的力。

硬约束的改变采用治理事件而非梯度。提案 $\Delta G_h$ 需要提出者、授权依据、影响范围、生效时间和撤销路径；承诺更新还要求相关主体确认。若“自由”和“安全”冲突，系统不能凭借两个矩阵不对易便宣称得到解，而要列出受影响主体、不可接受后果、可补偿方案和未解决分歧。多主体没有共同标尺时，输出应保留分歧结构。

我们把冲突诊断分为三层。第一层是内部不一致：同一版本规则对同一条件给出相反资格；第二层是跨版本漂移：旧承诺在新规则下被悄然删除；第三层是环境冲突：合法行动在当前资源或物理条件下不可执行。对应恢复分别是规则消歧、迁移审计和重新规划。焦虑、犹豫或高激活可能是人的经验现象，但不是软件系统中规范冲突成立的必要证据。

规范更新还需要反操纵边界。一次奖励、用户的临时命令或高频点击不能自动覆盖安全政策和长期承诺；模型生成的自我评价也不能充当独立授权。Boundary区分事实输入、偏好信号、授权指令和攻击载荷，TCE只能在当前资格集合内分配控制。这样，“目的驱动”就不会退化为谁提供更强信号谁接管系统。

经验验证至少包括偏好预测、干预响应和长期一致性。预测准确只能说明模型在样本分布中拟合反馈；要判断某条偏好是否因经历改变，需要控制混杂并比较反事实；要判断规范更新是否成功，还要检查承诺保持、受影响主体和异常条件。由此，原章关于行为反塑价值、价值冲突和慢尺度变化的问题被完整保留，但结论不再依赖杨—米尔斯同构。

AgentOS可把 $g$ 的稳定部分放入 $\Theta$，把事件性反馈放入 $E$，由SPC提供主体相关视图，由Boundary和权限系统执行硬约束，由外部记忆保存承诺证据。这是一套工程实例。其他系统可以使用制度规则、神经调节或群体投票，只要更新对象、授权来源、冲突程序和验证证据是明确的，就能映射到HSF-HD的一般规范动力学。

规范模型还要允许“未知”和“未决”。当证据不足以判断偏好、两个授权来源冲突或相关主体尚未响应时，系统不能用默认权重悄然制造一致意见。未决项进入带截止时间和责任人的队列，低风险行动可在可逆范围内继续，高风险行动暂停。这个机制把不确定性显式纳入动力学，也防止系统把沉默误读为同意。

长期一致性不等于永不改变。合法的新承诺、环境变化和主体反思都可以修改 $g$，关键是变化路径可解释、受权且保留旧版本。评审时比较的不只是前后参数距离，还包括哪些主体受益或受损、哪些例外被新增、哪些旧案例改变判定。这样才能区分真正的规范学习、策略漂移和外部操纵。

\section{主动控制算子：快慢回路、预算与行动资格}
\label{sec:section_9657}

控制问题是在给定状态估计、目标和资格集合下选择干预，而不是从状态与价值的换位子直接导出“意志”。设离散系统为
\begin{equation}
x_{k+1}=F_kx_k+B_ku_k+G_ke_k+\xi_k,\qquad u_k\in\mathcal U_k,
\end{equation}
其中 $u_k$ 具有明确通道和单位，$\mathcal U_k$ 汇总幅度、速率、权限和资源限制。有限时域控制器求解
\begin{equation}
u^*_{k:k+H-1}=\arg\min_{u\in\mathcal U_k}J_k(u)\quad
\text{s.t. }x_{j+1}=F_jx_j+B_ju_j+G_je_j,;x_j\in\mathcal X_{safe}.
\end{equation}
若模型不确定，约束应对置信集合或风险分位数成立；若求解超时，系统执行经过验证的后备策略，而不是使用过期最优解。

快回路改变活动增益、候选抑制和路由偏置。增益不是在流形上挖势阱，而是调整相关证据的读出权重；抑制不是建造势垒，而是在冲突或风险条件下限制传播与行动；偏置则改变搜索优先序。三者都必须保持校准：提高某候选的计算预算不能被报告成提高其真实性。对戒烟、紧急报警或工具选择等例子，系统分别需要目标证据、风险规则和权限，而不能只比较激活强度。

慢回路提出参数训练、结构编辑、外部化和接口重构。它不在连续积分步中直接写 $\dot S$，而是形成提案、影子验证和原子提交。Store与Update对应证据写入与版本修订，Create与Integrate对应新节点、关系或抽象接口，Boundary和Offloading决定信息是否跨边界、是否交给外部存储或专业工具。每个操作保存前置条件、影响集合和回滚指针。

控制预算至少分为计算时间、内存、外部调用、人工注意和设备能耗。前三者可以由软件计数，人工注意需要实验或流程估计，焦耳能耗必须由硬件测量或经校准的功率模型得到。它们可以相关，但不能因都被称为“成本”而相互替代。CCM根据延迟、候选数、未决验证和风险分配预算；资源不足时降低自动化等级、缩小任务或请求援助。

闭环控制要求独立反馈。行动 $u_k$ 执行后，系统读取 $o_{k+1}$，比较预期效果与实测结果，更新状态估计并判断提交、补偿或停止。若行动改变了环境，新的输入分布也随之改变；因此不能把环境只当外部源项。对不可逆操作，补偿可能只能减轻后果，日志应明确“内部已回滚、外部未复原”。

稳定性结论要附条件。若存在 $V(x)\ge0$ 使闭环期望满足
\begin{equation}
\mathbb E[V(x_{k+1})\mid x_k]-V(x_k)\le-\alpha\|x_k\|^2+\beta\|e_k\|^2,
\end{equation}
则在给定扰动界内可讨论均方有界；这不证明任务完成、规范合规或语义正确。任务层、状态层、规范层和资源层各有独立验收。所谓“意志”在本章只指跨这些层协调干预的控制能力，不是可由一条物理方程推出的实体。

在AgentOS中，TCE可实现滚动控制，SPC提供压缩观测，CCM约束搜索复杂度，Boundary决定行动资格，Offloading安排外部计算，$E$与$M_e$形成反馈证据。其他系统可能没有这些名称，但只要具有状态估计、目标、约束、控制、反馈和恢复，就属于同一控制模型族。

控制器自身也会失效。状态估计可能偏置，执行器可能饱和，模型服务可能超时，权限可能在计划期间撤销。每条控制命令因此带有效期、幂等键和执行前复核；执行器返回实际完成量，而不是简单成功标志。部分完成时，验证器根据外部状态选择续执行、补偿或停止，不能假设命令原子完成。

快慢回路之间还存在次序问题。先提高增益再写入记忆，与先更新结构再提高增益可能产生不同结果。系统通过影子运行估计操作交换差异，超过阈值时串行化，并把次序写入事件日志。对并发工具调用，版本向量和冲突合并规则承担同样作用。控制理论在这里与事务协议相接，防止连续方程掩盖实际竞态。

\section{混合闭环：联合算子、全耦合拓扑与验收边界}
\label{sec:section_7908}

把前述部分合并，系统不是四条物理场方程，而是一个连续流与离散事件交替的混合动力系统。固定版本区间内，活动、分布表示和连续参数按
\begin{equation}
\dot X_c=F(X_c,S,g,o,c)
\end{equation}
演化；当验证、权限、结构提案或环境回执触发事件 $e_j$ 时，系统执行
\begin{equation}
(X_c^+,S^+,g^+,M_e^+)=R_j(X_c^-,S^-,g^-,M_e^-,e_j).
\end{equation}
$R_j$ 是带前置条件的重置映射。连续模型负责局部预测，重置映射负责结构、规范和记忆版本；二者不能互相冒充。

一次完整循环包含观测、估计、候选生成、资格筛选、控制、执行、验证与提交。环境通过传感和工具结果影响证据，系统通过行动改变环境；活动状态影响训练样本选择，结构又改变后续传播；行为结果可以形成偏好证据，但规范变更需要授权；控制分配资源并受硬约束限制。原章所谓现实、结构、价值和控制四组双向耦合由此得到可执行解释。

\begin{figure}[htbp]
    \centering
    \resizebox{0.8\textwidth}{!}{\documentclass[tikz,border=10pt]{standalone}
\usepackage{tikz}
\usepackage{amsmath}
\usepackage{amssymb}
\usepackage{ctex} % 中文支持

\begin{document}
% 加载必要的 TikZ 库
\usetikzlibrary{shapes.geometric, arrows.meta, positioning, calc, shadows, backgrounds, fit}

% --- 颜色定义 (HSF-HD 风格) ---
\definecolor{hsfBlue}{RGB}{22, 160, 133}   % 几何/结构
\definecolor{hsfRed}{RGB}{192, 57, 43}     % 外部/激波
\definecolor{hsfPurple}{RGB}{142, 68, 173} % 价值/目的
\definecolor{hsfOrange}{RGB}{211, 84, 0}   % 意志/控制
\definecolor{hsfDark}{RGB}{44, 62, 80}     % 核心思维



\begin{tikzpicture}[
    node distance=2.5cm and 4cm,
    % --- 节点样式 ---
    core/.style={
        circle, 
        draw=hsfDark, 
        ultra thick, 
        fill=hsfDark!10, 
        minimum size=3.5cm, 
        align=center, 
        drop shadow,
        font=\bfseries
    },
    component/.style={
        rectangle, 
        rounded corners=8pt, 
        draw=black!60, 
        thick, 
        minimum width=4cm, 
        minimum height=1.8cm, 
        align=center, 
        drop shadow,
        font=\small
    },
    % --- 连线样式 ---
    link/.style={
        <->, 
        >={Latex[length=3mm, width=2mm]}, 
        line width=1.5pt, 
        draw=gray!80
    },
    dashed_link/.style={
        ->, 
        >={Latex[length=2mm, width=1.5mm]}, 
        line width=1pt, 
        draw=gray!60,
        dashed
    },
    % --- 标签样式 ---
    label_text/.style={
        fill=white, 
        fill opacity=0.9, 
        text opacity=1, 
        font=\footnotesize\sffamily, 
        align=center,
        inner sep=2pt,
        text=black!80
    }
]

    % =================================================
    % 1. 核心节点：TDE (思维流)
    % =================================================
    \node[core] (TDE) {
        \textcolor{hsfDark}{\Large \textbf{TDE}}\\
        \textbf{目的论狄拉克方程}\\
        \textit{思维流 $\Psi$}\\
        (推理/演化)
    };

    % =================================================
    % 2. 周围节点：内部组件
    % =================================================
    
    % 左侧：CEFE (结构)
    \node[component, left=of TDE, fill=hsfBlue!10, draw=hsfBlue] (CEFE) {
        \textcolor{hsfBlue}{\large \textbf{CEFE}}\\
        \textbf{认知爱因斯坦方程}\\
        \textit{度量 $g_{\mu\nu}$}\\
        (记忆/习惯/形)
    };

    % 右侧：CYME (价值)
    \node[component, right=of TDE, fill=hsfPurple!10, draw=hsfPurple] (CYME) {
        \textcolor{hsfPurple}{\large \textbf{CYME}}\\
        \textbf{认知杨-米尔斯方程}\\
        \textit{规范势 $\mathcal{A}_\mu$}\\
        (偏好/目的/质)
    };

    % 下方：TCE (意志)
    \node[component, below=of TDE, fill=hsfOrange!10, draw=hsfOrange, yshift=-0.5cm] (TCE) {
        \textcolor{hsfOrange}{\large \textbf{TCE}}\\
        \textbf{目的论控制方程}\\
        \textit{宏观算子 $\hat{\mathcal{O}}$}\\
        (控制/注意/元认知)
    };

    % =================================================
    % 3. 外部节点：环境
    % =================================================
    
    % 上方：ENV (环境)
    \node[component, above=of TDE, fill=hsfRed!10, draw=hsfRed, yshift=0.5cm] (ENV) {
        \textcolor{hsfRed}{\large \textbf{External Env}}\\
        \textbf{外部物理世界}\\
        \textit{源流 $\vec{J}_{ext}$ / 动作 Action}
    };

    % =================================================
    % 4. 绘制背景框 (区分内外)
    % =================================================
    \begin{scope}[on background layer]
        \node[fit=(CEFE)(CYME)(TCE)(TDE), 
              draw=gray!30, dashed, rounded corners=15pt, 
              fill=gray!5, inner sep=15pt,
              label={[anchor=north west, xshift=10pt, yshift=-5pt, text=gray]north west:\textbf{智能体内部 (Internal Manifold)}}] (box) {};
    \end{scope}

    % =================================================
    % 5. 绘制连线 (主要耦合) - 星型拓扑无交叉
    % =================================================

    % I. 现实耦合 (ENV <-> TDE)
    \draw[link, draw=hsfRed] (ENV) -- (TDE) 
        node[midway, label_text] {I. 现实耦合\\(感知 $\downarrow$ / 行动 $\uparrow$)};

    % II. 结构耦合 (TDE <-> CEFE)
    \draw[link, draw=hsfBlue] (TDE) -- (CEFE) 
        node[midway, label_text] {II. 结构耦合\\(认知应力 $\to$ / $\leftarrow$ 几何惯性)};

    % III. 价值耦合 (TDE <-> CYME)
    \draw[link, draw=hsfPurple] (TDE) -- (CYME) 
        node[midway, label_text] {III. 价值耦合\\(价值流 $\to$ / $\leftarrow$ 洛伦兹力)};

    % IV. 控制耦合 (TDE <-> TCE)
    \draw[link, draw=hsfOrange] (TDE) -- (TCE) 
        node[midway, label_text] {IV. 控制耦合\\(偏差检测 $\downarrow$ / $\uparrow$ 算子注入)};

    % =================================================
    % 6. 隐式/次级耦合 (可选，使用曲线避免交叉)
    % =================================================
    
    % TCE -> CEFE (塑性重构)
    \draw[dashed_link, bend left=45] (TCE.west) to node[label_text, font=\tiny] {V. 塑性重构} (CEFE.south);
    
    % TCE -> CYME (价值重估)
    \draw[dashed_link, bend right=45] (TCE.east) to node[label_text, font=\tiny] {VI. 价值重估} (CYME.south);

\end{tikzpicture}

\end{document}}
    \caption{全耦合拓扑：图保留原有五类接口，正文将其解释为观测、状态、结构、规范与控制之间的可验证信息流，而非物理场同构}
\end{figure}

图中的中心状态不是单一“思维实体”。局部结构网络与全局分布表示通过身份绑定和读出联合，环境输入携带来源与时序，控制输出携带权限与补偿策略。任何箭头都对应接口契约：输入类型、版本、最大延迟、误差界和失败处理。没有这些契约，拓扑图只能说明概念关联，不能证明闭环能够运行。

多时间尺度的充分分离也不是先验事实。若快状态在固定慢变量下没有稳定吸引集，或结构更新与活动更新频率接近，冻结慢变量的近似会失效；若反馈延迟超过闭环裕量，控制可能振荡；若规范频繁改变，历史评价集会失去可比性。因此部署前要通过系统辨识估计时标、延迟与耦合强度，并在超出适用域时重新建模。

以长期研究任务为例，快速回路读取论文和生成假设，中速回路维护问题树、来源绑定与实验计划，慢速回路在复现实验后修改稳定知识和方法规范。外部工具故障可能造成观测缺口，新的反例可能触发结构重组，伦理或预算变化可能收缩行动集合。系统即使找到更低损失的解释，也只有在来源、反例、权限和实验结果均通过时才能提交“任务成功”。这个例子保留了原章从感知到行动、从活动到记忆、从行为到价值、从意图到控制的超循环，但每个环节都有可审计证据。

验收分为五层：数值层检查步长、溢出和随机性；状态层检查校准、绑定与有界性；结构层检查版本迁移和保留查询；规范层检查权限、承诺与硬约束；任务层检查外部结果。通过其中一层不能替代其他层。预测有效只支持模型在给定任务、数据和干预范围内可用，不能推出整个智能系统与状态方程同构。

因此，本章的统一不是把所有机制压成一个“上帝公式”，而是给出共同的状态语法和组合规则：对象先分型，目标与约束先声明，连续变化与离散事件分开，快慢回路通过验证边界交换信息，局部结构与全局表示通过绑定接口联合，控制在资格和预算内执行。HSF-HD作为智能的一般“空气动力学”提供这些可迁移关系；AgentOS把它们实现为具体运行组件。不同载体只要能报告对应状态、接口、干预和边界，就可以比较，而不必共享同一物理隐喻。

闭环还要接受故障注入。我们分别中断传感输入、延迟外部记忆、污染单一来源、撤销权限、制造结构迁移失败和耗尽调用预算，观察系统是否进入预定降级状态。若只有正常路径有效，就不能声称动力学内核稳健。每种故障都应对应可观测告警、最大恢复时间和数据一致性检查，恢复后还要重放未决事务。

最终可迁移的不是某组符号，而是一套提问顺序：状态是什么，谁能观测，哪些更新连续，哪些更新是事件，目标和硬约束由谁授权，误差怎样传播，何时提交，失败怎样恢复。任何实现只要回答这些问题，就能与本章模型比较；回答不了时，再相似的场论形式也只是隐喻。这个边界保持了理论的广度，同时让每个局部结论都能被实验、日志和反例检验。
